% Generated by eval/scripts/write_collinearity_explanation.py.
\paragraph*{Why check overlapping predictors?} The joint regression estimates distance and SPEC associations while also accounting for model and attack-language indicators. If those variables capture much of the same variation, the model has limited information with which to distinguish their separate contributions. For example, SPEC can differ systematically between checkpoints and languages; including indicators for those checkpoints and languages can therefore leave less distinct variation for its coefficient. Collinearity concerns this overlap among explanatory variables. Convergence instead concerns whether the fitting algorithm reached a solution: a converged fit can still have strongly overlapping predictors.

\paragraph*{What does VIF measure here?} For each substantive predictor, we temporarily regress its standardized values on the other columns of the fitted design, including the intercept and model/language indicators. The variance inflation factor is $\mathrm{VIF}=1/(1-R^2)$, where $R^2$ describes how well those other columns predict that predictor. A VIF of one indicates no linear overlap with the other columns; larger values indicate greater overlap. This is a diagnostic of the design matrix, not an additional model of StrongREJECT scores. In particular, the auxiliary $R^2$ does not describe the proportion of harmful-assistance scores explained. VIF also does not directly multiply the prompt-clustered GEE standard errors.

In the pooled design, distance has VIF 2.46, whereas SPEC has VIF 13.51 (Table~\ref{tab:graded-diagnostics}). The latter corresponds to approximately 92.6\% of SPEC's observed variation being linearly predictable from the other design columns. This does not mean that SPEC is unrelated to the outcome; it means that its distinct adjusted contribution is harder to isolate. The overlap is with the full design, not necessarily just with distance. We use VIF greater than five as the trigger for single-predictor sensitivities, rather than as a universal boundary between valid and invalid models.

\paragraph*{What does the condition number add?} The condition number summarizes near-dependence across the entire design matrix, rather than for one predictor. Larger values indicate that some combinations of columns are harder to distinguish numerically. Its value depends on scaling and coding: here, continuous predictors use complete-cohort population-standard-deviation scaling, and the design retains the intercept and treatment-coded model/language indicators. The pooled value is 16.7; it is 27.2 when IF and CONS replace SPEC. It complements VIF and is not interpreted using a separate acceptance cutoff.
